\section{Results}
\label{sec:results}

\subsection{Super-Resolution}

All 9 models (Table~\ref{tab:experiments}) are evaluated on the test set.
Metrics: PSNR, SSIM, LPIPS~\cite{zhang2018}, and MSE.
Table~\ref{tab:metrics} presents the full results.

\begin{table}[H]
\centering
\caption{Quantitative evaluation of all 9 SR models on the test set (Phase 1, L1 loss).}
\label{tab:metrics}
\footnotesize
\renewcommand{\arraystretch}{1.0}
\begin{tabular}{@{}lcccc@{}}
\toprule
\textbf{Model} & \textbf{PSNR (dB)} $\uparrow$ & \textbf{SSIM} $\uparrow$ & \textbf{LPIPS} $\downarrow$ & \textbf{MSE} $\downarrow$ \\
\midrule
M1 & 35.19 & 0.922 & 0.0601 & 3.18e-4 \\
M2 & 30.84 & 0.817 & 0.2064 & 9.12e-4 \\
M3 & 26.93 & 0.666 & 0.4887 & 2.58e-3 \\
\midrule
M4 & \textbf{37.43} & \textbf{0.944} & \textbf{0.0429} & \textbf{1.91e-4} \\
M5 & 32.01 & 0.860 & 0.1379 & 6.96e-4 \\
M6 & 27.49 & 0.729 & 0.4135 & 2.19e-3 \\
\midrule
M7 & 34.79 & 0.918 & 0.0799 & 3.48e-4 \\
M8 & 31.20 & 0.855 & 0.1709 & 8.16e-4 \\
M9 & 27.52 & 0.765 & 0.4100 & 2.09e-3 \\
\bottomrule
\end{tabular}
\end{table}

\textbf{Best model: M4} (Scenario B, $\times$2).
PSNR = 37.43\,dB, SSIM = 0.944, LPIPS = 0.0429, MSE = 1.91e-4.

\subsubsection{Scale Factor Impact}

\begin{itemize}
    \item $\times$\textbf{2} (avg): PSNR 35.80\,dB, SSIM 0.928, LPIPS 0.061 --- \textit{excellent}
    \item $\times$\textbf{4} (avg): PSNR 31.35\,dB, SSIM 0.844, LPIPS 0.172 --- \textit{good}
    \item $\times$\textbf{8} (avg): PSNR 27.31\,dB, SSIM 0.720, LPIPS 0.437 --- \textit{limited}
\end{itemize}

$\times$2 to $\times$4 drops 4.45\,dB. $\times$4 to $\times$8 drops 4.04\,dB.
Higher scale factors are increasingly ill-posed.

\subsubsection{Scenario Comparison}

Scenario B (bilinear) achieves the highest average PSNR (32.31\,dB) and lowest LPIPS (0.198).
Scenario A (bicubic): 30.99\,dB. Scenario C (Lanczos): 31.17\,dB.
Bilinear introduces the least aliasing, producing smoother LR inputs that are easier to reconstruct.
Scenario C maintains competitive SSIM (0.846 avg) despite the most aggressive degradation.

\subsection{Visual Results}

\begin{figure}[H]
\centering
\includegraphics[width=\columnwidth]{sr_comparison.png}
\caption{LR $\rightarrow$ SR $\rightarrow$ HR for best ($\times$2, top) and worst ($\times$8, bottom) models. Zoom patches show fine detail recovery.}
\label{fig:sr_results}
\end{figure}

At $\times$2, SR outputs closely approximate HR reference.
Vessel structures and optic disc boundaries are visibly recovered.
At $\times$8, fine vessels are not fully restored, but overall fundus structure and lesion contrast are preserved.

\subsection{DR Classification}

ResNet50 ordinal regression achieves 73.56\% accuracy and QWK = 0.8764 on the test set (537 images).
This improves over the VGG16 baseline (64.43\%) by +9.13\%.

\begin{table}[H]
\centering
\caption{ResNet50 ordinal regression per-class results on the test set.}
\label{tab:classification}
\footnotesize
\renewcommand{\arraystretch}{1.0}
\begin{tabular}{@{}lcccc@{}}
\toprule
\textbf{Class} & \textbf{Support} & \textbf{Precision} & \textbf{Recall} & \textbf{F1-score} \\
\midrule
No DR (0)        & 271 & 0.960 & 0.978 & 0.969 \\
Mild (1)         &  53 & 0.518 & 0.547 & 0.532 \\
Moderate (2)     & 143 & 0.728 & 0.525 & 0.610 \\
Severe (3)       &  27 & 0.232 & 0.815 & 0.361 \\
Proliferative (4)&  43 & 0.571 & 0.093 & 0.160 \\
\midrule
Weighted avg     & 537 & 0.787 & 0.736 & 0.735 \\
\bottomrule
\end{tabular}
\end{table}

Binary DR detection (classes 1--4 vs.~0) reaches 96.83\% accuracy, 95.86\% sensitivity, and 97.79\% specificity.
Class 0 (no DR) is reliably classified (F1 = 0.969).
Minority classes (mild, severe, proliferative) show lower F1-scores due to limited training samples (5.1\%--9.8\%).
